\documentclass{article}
\usepackage{import}
\usepackage{tcolorbox}
\usepackage{amsmath}
\usepackage{amssymb}
\usepackage{amsthm}
\usepackage{geometry}
\usepackage{enumitem}
\usepackage{pdfpages}
\usepackage{transparent}
\usepackage{booktabs}
\usepackage{bm}
\usepackage{tabularx}
\usepackage{multirow}
% \usepackage{minted}
\usepackage{xcolor}
\usepackage{setspace}
\usepackage[hidelinks]{hyperref}
\usepackage{pgfplots}
% \usepackage{soulpos}

\tcbuselibrary{theorems}
\tcbuselibrary{skins}

\setcounter{tocdepth}{6}
\setcounter{secnumdepth}{6}

\newcommand{\param}{\boldsymbol{\theta}}
\newcommand{\bs}{\boldsymbol}
\newcommand{\data}{\mathcal D}

% pgfplots settings
\pgfplotsset{compat=1.18}
\usepgfplotslibrary{external}

% tabularx Settings
\renewcommand\tabularxcolumn[1]{>{\centering\arraybackslash}m{#1}}

% Page Settings
\newgeometry{vmargin={15mm}, hmargin={15mm, 15mm}}
\setlength{\parindent}{0pt}
\setlength{\parskip}{0.7\baselineskip}%{0.8em plus 1pt minus 1pt}

\onehalfspacing

% hyperref Settings
\hypersetup{
    colorlinks=true,
    linkcolor=blue,
    filecolor=magenta,      
    urlcolor=cyan,
}

% inkscape Settings
\newcommand{\incfig}[1]{%
    \def\svgwidth{\columnwidth}
    \import{./figure/}{#1.pdf_tex}
}

% inline code block settings
\newtcbox{\mybox}[1][]{
    on line,
    arc=1pt, outer arc=2pt,
    colback=lightgray!50!white, colframe=lightgray!50!white,
    boxsep=0pt, left=0pt, right=-1pt, top=2pt, bottom=0pt,
    boxrule=0pt, #1
}

\NewDocumentCommand{\code}{v}{%
  \begingroup\ttfamily
  \codeulinner{#1}%
  \endgroup%
}

% Definition Box Styles
\newtcbtheorem[auto counter, number within=subsection]{defn}{Definition}{
% basic settings
    theorem name and number,
    separator sign={},
    description delimiters={(}{)},
    enhanced jigsaw,
    sharp corners,
% box margins and rules
    boxrule=1pt,
    left=0.2cm,right=0.2cm,top=0.2cm,
    toptitle=0.1cm,
    bottomtitle=0.08cm,
% box colors
    colframe=white!25!black,
    colback=white,
    coltitle=white,
% box fonts
    fonttitle=\bfseries,fontupper=\normalsize
}{defn}

% Theorem Box Styles
\newtcbtheorem[auto counter, number within=subsection]{thm}{Theorem}{
% basic settings
    theorem name and number,
    separator sign={},
    description delimiters={(}{)},
    enhanced jigsaw,
    sharp corners,
% box margins and rules
    boxrule=1pt,
    left=0.2cm,right=0.2cm,top=0.2cm,
    toptitle=0.1cm,
    bottomtitle=0.08cm,
% box colors
    colframe=white!25!black,
    colback=white,
    coltitle=white,
% box fonts
    fonttitle=\bfseries,fontupper=\normalsize
}{thm}

% Lemma Box Styles
\newtcbtheorem[number within=tcb@cnt@thm]{lem}{Lemma}{
% basic settings
    theorem name and number,
    separator sign={},
    description delimiters={(}{)},
    enhanced jigsaw,
    sharp corners,
% box margins and rules
    boxrule=1pt,
    left=0.2cm,right=0.2cm,top=0.2cm,
    toptitle=0.1cm,
    bottomtitle=0.08cm,
% box colors
    colframe=white!50!black,
    colback=white,
    coltitle=white,
% box fonts
    fonttitle=\bfseries,fontupper=\normalsize
}{lem}

% Corollary Box Styles
\newtcbtheorem[number within=tcb@cnt@thm]{cor}{Corollary}{
% basic settings    
    theorem name and number,
    separator sign={},
    description delimiters={(}{)},
    enhanced jigsaw,
    sharp corners,
% box margins and rules
    boxrule=1pt,
    left=0.2cm,right=0.2cm,top=0.2cm,
    toptitle=0.1cm,
    bottomtitle=0.08cm,
% box colors
    colframe=white!50!black,
    colback=white,
    coltitle=white,
% box fonts
    fonttitle=\bfseries,fontupper=\normalsize
}{cor}

% Remark Box Styles
\newtcbtheorem[number within=subsection]{rem}{Remark}{
% basic settings
    theorem name and number,
    separator sign={},
    description delimiters={(}{)},
    enhanced jigsaw,
    sharp corners,
% box colors
    colframe=white!60!black,
    colback=white,
    coltitle=black,
% box margins and rules
    boxrule=1pt,
    left=0.2cm,right=0.2cm,top=0.2cm,
    toptitle=0.1cm,
    bottomtitle=0.08cm,
% box fonts
    fonttitle=\bfseries,fontupper=\normalsize
}{rem}

% Example Box Styles
% Remark Box Styles
\newtcbtheorem[number within=subsection]{exm}{Example}{
% basic settings
    theorem name and number,
    separator sign={},
    description delimiters={(}{)},
    enhanced jigsaw,
    sharp corners,
% box colors
    colframe=white!60!black,
    colback=white,
    coltitle=black,
% box margins and rules
    boxrule=1pt,
    left=0.2cm,right=0.2cm,top=0.2cm,
    toptitle=0.1cm,
    bottomtitle=0.08cm,
% box fonts
    fonttitle=\bfseries,fontupper=\normalsize
}{exm}


% Proof Box Styles

\title{Normalization Methods}
\author{Hyoung Joon, Kim}
\date{\today}

\begin{document}
\maketitle
\tableofcontents
\newpage

\section{Introduction}
    We introduce methods of \textbf{normalization}, which helps deep neural network to be trained.
    Three methods will be discussed: \textbf{Data normalization}, 
    \textbf{layer normalization} and \textbf{batch normalization}.

\section{Main Ideas}
    \subsection{Data Normalization}
        \textbf{Data normalization} is to normalize the test data along mini-batch dataset.
        Since each feature values are distributed along different ranges,
        some feature may be ignored, or overemphasized. Normalizing the data solves this probelm, by
        making all the feature values to have zero mean and unit variance.
        
        \emph{My experience about the effect of data normalization:
        Polynomial regression exploded at degree 7 without data normalization.
        But with data normalization, it did not explode even when I used degree 500 for polynomial features.}

    \subsection{Batch Normalization}\label{sec:batchNorm}
        \textbf{Batch normalization} is to normalize the hidden layers about the batch dataset.
        We first forward pass the mini-batch data, and then, we calculate the mean and standard
        deviation of the hidden units calculated about each data in the mini-batch dataset.
        We can apply this at pre-activation or activated layer. In practice, either method can be used.

        Let $a_ni$ be the $i$-th pre-activation unit calculated with $n$-th element of mini-batch dataset.
        Then,
        \begin{align}
            \mu_i &= \frac{1}{K}\sum_{n=1}^{K}a_{ni} \\
            \sigma_i^2 &= \frac{1}{K}\sum_{n=1}^{K}(a_{ni}-\mu_i)^2
        \end{align}
        After that, we can calculate the normalized units: 
        \[
        \hat{a}_{ni} = \frac{a_{ni}-\mu_i}{\sigma_i + \epsilon}
        \]
        The normalized unit losts some of the degree of freedom. To compensate this, we can scale item
        by using two learnable parameters $\gamma_i$ and $\beta_i$;
        \[
        \tilde{a}_{ni} = \gamma_i \hat{a}_{ni} + \beta_i
        \]
        \subsubsection{Inference}
            In the inference step, we cannot use batch normalization, since there only a single data is
            used. To solve this, we can evaluate the mean and stdev across the whole hidden units.
            However, this is usually expensive, so we use \textbf{moving averages} calculated during the
            training process:
            \begin{align}
                \bar{\mu}^{(\tau)}_i &= \alpha\bar{\mu}^{(\tau-1)}_i + (1-\alpha)\mu_i \\
                \bar{\sigma}^{(\tau)}_i &= \alpha\bar{\sigma}^{(\tau-1)}_i + (1-\alpha)\sigma_i
            \end{align}
            calculated with $\tau$-th mini-batch dataset, where $0 \leq \alpha \leq 1$.
        
    \subsection{Layer Normalization}
            Using batch normalization injects too much noise if the batch size is too small.
            Even when the dataset is too large, the data can be distributed to different GPUs, and
            batch normalization cannot be applied to the layer. \textbf{Layer normalization} is
            the alternative of batch normalization, which also uses training datasets to normalize
            the weights and biases.

            Layer normalization normalize the hidden units with respect to the layer itself.
            Let's assume that we have hidden layer with $M$ hidden units. Then, we calculate
            the mean and variance of hidden units;
            \begin{align}
                \mu_n&=\frac{1}{M}\sum_{i=1}^{M}a_{ni} \\
                \sigma_n^2&=\frac{1}{M}\sum_{i=1}^{M}(a_{ni}-\mu_i)^2
            \end{align}
            Then, as we did in batch normalization, we normalize the hidden units;
            \[
            \hat{a}_{ni} = \frac{a_{ni}-\mu_n}{\sigma_n + \epsilon}
            \]
            Note that the layer normalization an be applied during inference,
            since we can apply layer normalization to one data.
\end{document}